%% kapa: exam=1 round=22 year=2011 period=2 form=A subject=accounting subject_ko=회계학 seq=6 number=6
%% source: Q-Net 감정평가사 기출문제 (cst00309) artlSeq=5036790 file=extracted/1008/5036790/2011년 제22회 감정평가사 1차시험 기출문제/2교시.hwp
%% conversion: kordoc 4.24.51 -> tools/gen1.py (auto); review=proofread-gpt-astra; flags=none
\begin{question}{6}
다음 자료를 기초로 (주)서울의 포괄손익계산서에 계상될 퇴직급여비용을 계산하면?
\begin{framedtext}
ㆍ당기근무원가 ￦450,000 \par
ㆍ기초 확정급여채무 금액 ￦900,000 (할인율 연 10\%) \par
ㆍ사외적립자산의 실제수익 ￦150,000, 기대수익 ￦120,000 \par
ㆍ회사는 보험수리적손익을 전액 기타포괄손익으로 인식하고 있음 \par
ㆍ과거 근무원가 중 당기손익인식금액 ￦80,000 \par
\end{framedtext}
\begin{choices}
  \item ￦424,000
  \item ￦444,000
  \item ￦450,000
  \item ￦494,200
  \item ￦500,000
\end{choices}
\end{question}
\answer{5}
